% ============================================================
% CHAPITRE 3 — Sprint 0 : Étude préliminaire
% ============================================================
\chapter{Sprint 0 --- Étude Préliminaire}

\section{Introduction}

Le Sprint~0 est une étape fondamentale dans la planification et l'organisation
du développement de la plateforme de gestion du recrutement DAAM. Ce sprint
ne produit pas encore de fonctionnalités opérationnelles livrables aux utilisateurs
finaux, mais sert à établir les bases du projet~: identification des besoins,
définition des fonctionnalités clés, choix technologiques, ainsi que la
planification des itérations.

Cette phase permet également de construire une vision globale du système à
travers des modèles de cas d'utilisation et un Product Backlog structuré. Elle
prépare ainsi les sprints suivants, consacrés à la réalisation progressive des
modules métier.


\section{Analyse des besoins}

L'analyse des besoins est essentielle pour assurer que la plateforme réponde
aux attentes des utilisateurs finaux et au contexte multi-agences de DAAM.
Elle consiste à identifier les différents acteurs du système, puis à préciser
leurs besoins fonctionnels et non fonctionnels.


\subsection{Identification des acteurs}

Trois types d'utilisateurs principaux ont été identifiés dans le système~:

\begin{itemize}
  \item \textbf{Administrateur}
  \begin{itemize}
    \item Dispose des droits de gouvernance globale de la plateforme.
    \item Gère les utilisateurs, les zones, les agences et les affectations RH.
    \item Consulte une vue transverse sur les recrutements, les candidatures
          et les candidats admis, organisée par responsable RH.
    \item Accède au dashboard global de sélection.
  \end{itemize}

  \item \textbf{RH (Responsable Ressources Humaines)}
  \begin{itemize}
    \item Gère les opérations quotidiennes de recrutement sur son périmètre
          (zones / agences affectées).
    \item Crée et publie les offres, traite les candidatures, planifie les
          entretiens et confirme les embauches.
    \item Utilise le dashboard, le calendrier, les QCM et le suivi des
          candidats admis.
  \end{itemize}

  \item \textbf{Candidat}
  \begin{itemize}
    \item Consulte les offres publiées.
    \item Dépose sa candidature avec CV et informations personnelles.
    \item Suit l'état de ses dossiers et participe aux évaluations / QCM.
  \end{itemize}
\end{itemize}


\subsection{Définition des besoins fonctionnels}

Les besoins fonctionnels décrivent les fonctionnalités essentielles que la
plateforme doit offrir. Ils couvrent l'ensemble du cycle de vie du recrutement,
depuis l'authentification jusqu'au suivi post-embauche.


\subsubsection{Gestion des comptes et de l'authentification}

\begin{itemize}
  \item \textbf{Authentification sécurisée}~: les utilisateurs doivent pouvoir
        s'authentifier de manière sécurisée (approche basée sur JWT).
  \item \textbf{Gestion de profil}~: chaque utilisateur doit pouvoir consulter
        et mettre à jour ses informations personnelles.
  \item \textbf{Administration des comptes}~: l'administrateur doit pouvoir
        créer, modifier et activer/désactiver des comptes selon les rôles.
  \item \textbf{Gestion des comptes candidats (côté RH)}~: le RH doit pouvoir
        consulter les comptes candidats et gérer leur statut (activation /
        désactivation) lorsqu'une action réglementaire ou opérationnelle
        l'exige.
\end{itemize}


\subsubsection{Gestion organisationnelle (zones, agences, affectations)}

\begin{itemize}
  \item \textbf{Gestion des zones}~: l'administrateur doit pouvoir créer et
        administrer les zones géographiques / organisationnelles.
  \item \textbf{Gestion des agences}~: chaque agence doit être rattachée à une
        zone, avec ses informations (adresse, localisation éventuelle).
  \item \textbf{Affectation RH}~: un responsable RH doit pouvoir être associé
        à une ou plusieurs zones afin de garantir le respect du périmètre
        d'accès.
\end{itemize}


\subsubsection{Gestion des recrutements (offres)}

\begin{itemize}
  \item \textbf{Création et édition d'offres}~: le RH doit pouvoir créer des
        offres détaillées (intitulé, description, compétences, disponibilité,
        langues, références, etc.).
  \item \textbf{Publication / clôture}~: une offre doit pouvoir passer d'un
        état brouillon à publié, puis être clôturée.
  \item \textbf{Rattachement organisationnel}~: chaque offre doit être liée
        à une agence / zone.
  \item \textbf{Association d'un QCM}~: une offre peut être liée à un
        questionnaire de présélection.
\end{itemize}


\subsubsection{Gestion des candidatures}

\begin{itemize}
  \item \textbf{Dépôt de candidature}~: le candidat doit pouvoir postuler à une
        offre publiée en fournissant son CV et ses réponses éventuelles au QCM.
  \item \textbf{Suivi des dossiers}~: le candidat doit pouvoir consulter le
        statut de ses candidatures.
  \item \textbf{Traitement RH}~: le RH doit pouvoir consulter, filtrer et
        analyser les candidatures (score QCM, analyse CV, historique).
  \item \textbf{Workflow de statuts}~: une candidature doit évoluer selon des
        états métier clairs (en attente, retenu, non retenu, désisté, admis).
\end{itemize}


\subsubsection{Entretiens, meetings et embauche}

\begin{itemize}
  \item \textbf{Planification d'entretiens}~: le RH doit pouvoir planifier un
        entretien en ligne ou physique, avec horaires et lien éventuel.
  \item \textbf{Calendrier}~: les entretiens doivent être visibles dans un
        calendrier RH.
  \item \textbf{Confirmation d'embauche}~: après acceptation et entretien
        physique le cas échéant, le RH doit pouvoir confirmer l'embauche
        (contrat, date de prise de fonction, etc.).
  \item \textbf{Suivi des admis}~: les candidats embauchés doivent rester
        accessibles pour le suivi post-recrutement (évaluations, rapports).
\end{itemize}


\subsubsection{QCM, évaluation et aide à la décision}

\begin{itemize}
  \item \textbf{Gestion des QCM}~: le RH doit pouvoir créer et maintenir des
        questionnaires.
  \item \textbf{Évaluations}~: le candidat doit pouvoir passer des évaluations
        associées à son parcours.
  \item \textbf{Analyse assistée du CV}~: le système doit pouvoir extraire et
        comparer des compétences pour aider à la présélection.
  \item \textbf{Tableaux de bord}~: Admin et RH doivent disposer d'indicateurs
        (volumes, retenus, non retenus, désistés, taux de sélection).
\end{itemize}


\subsubsection{Exports et reporting}

\begin{itemize}
  \item \textbf{Export Excel}~: le RH doit pouvoir exporter des suivis de
        candidatures / évaluations.
  \item \textbf{Historique candidat}~: Admin et RH doivent pouvoir consulter
        l'historique des candidatures d'un même profil.
\end{itemize}


\subsubsection{Perspectives (hors cœur recrutement initial)}

\begin{itemize}
  \item \textbf{Module de réclamations}~: livré au Sprint~5 (chapitre~8) sous
        forme de microservice dédié.
  \item \textbf{Industrialisation DevOps}~: Docker, Jenkins et SonarQube
        (chapitre~9).
  \item \textbf{Notifications in-app plus complètes}~: à enrichir au-delà des
        e-mails métier déjà présents (entretien, embauche).
\end{itemize}


\subsection{Définition des besoins non fonctionnels}

Ces besoins garantissent la qualité du système au-delà de ses fonctions de base~:

\begin{itemize}
  \item \textbf{Sécurité}~: authentification, contrôle d'accès par rôle et
        protection des données personnelles des candidats.
  \item \textbf{Confidentialité}~: isolation des données selon le périmètre RH
        (zones) et séparation claire Admin / RH / Candidat.
  \item \textbf{Utilisabilité}~: interfaces claires, parcours simple pour les
        candidats, navigation structurée pour RH et Admin.
  \item \textbf{Maintenabilité}~: architecture modulaire (frontend / microservices)
        facilitant l'évolution.
  \item \textbf{Performance}~: temps de réponse acceptables pour la consultation
        des listes, dashboards et dépôts de candidatures.
  \item \textbf{Évolutivité}~: possibilité d'ajouter des modules (réclamations,
        notifications) sans remettre en cause le cœur métier.
  \item \textbf{Traçabilité}~: conservation de l'historique des candidatures,
        décisions et entretiens.
\end{itemize}


\section{Diagramme de cas d'utilisation globale}

Le diagramme de cas d'utilisation globale fournit une vue d'ensemble des
interactions entre les acteurs et le système. Il synthétise les grandes
familles de fonctionnalités~: authentification, administration organisationnelle,
gestion des offres, candidatures, entretiens, embauche, dashboards et
évaluations.

\clearpage
\begin{figure}[p]
  \centering
  \reportfig{img/uc-global-hires.jpg}
  \caption{Diagramme de cas d'utilisation globale de la plateforme DAAM}
  \label{fig:uc-global}
\end{figure}
\clearpage


\section{Product Backlog}

Le Product Backlog est une liste priorisée des fonctionnalités à implémenter.
Il est structuré autour de User Stories, représentant les besoins exprimés par
les acteurs. La priorisation privilégie d'abord les fondations (auth, organisation),
puis le cœur métier (offres, candidatures), ensuite le pilotage (entretiens,
embauche, dashboards).

\begin{table}[htbp]
  \centering
  \small
  \caption{Product Backlog du projet (extrait priorisé)}
  \label{tab:product-backlog}
  \begin{tabular}{|c|p{3.4cm}|p{6.2cm}|c|}
    \hline
    \textbf{ID} & \textbf{Fonctionnalité} & \textbf{Description} & \textbf{Priorité} \\
    \hline
    US-01 & Authentification & En tant qu'utilisateur, je m'authentifie et j'accède à mon espace selon mon rôle. & Haute \\
    \hline
    US-02 & Profil & En tant qu'utilisateur, je consulte et mets à jour mon profil. & Haute \\
    \hline
    US-03 & Gestion utilisateurs & En tant qu'Admin, je gère les comptes et les rôles. & Haute \\
    \hline
    US-04 & Zones / agences & En tant qu'Admin, je gère zones et agences. & Haute \\
    \hline
    US-05 & Affectation RH & En tant qu'Admin, j'affecte un RH à des zones. & Haute \\
    \hline
    US-06 & Offres RH & En tant que RH, je crée, publie et clôture des offres. & Haute \\
    \hline
    US-07 & QCM & En tant que RH, je crée un QCM et l'associe à une offre. & Moyenne \\
    \hline
    US-08 & Candidature & En tant que candidat, je postule avec mon CV. & Haute \\
    \hline
    US-09 & Traitement candidatures & En tant que RH, je suis les statuts et l'historique. & Haute \\
    \hline
    US-10 & Entretiens & En tant que RH, je planifie un meeting / entretien. & Haute \\
    \hline
    US-11 & Embauche & En tant que RH, je confirme l'embauche d'un candidat. & Haute \\
    \hline
    US-12 & Candidats admis & En tant que RH/Admin, je suis les profils embauchés. & Moyenne \\
    \hline
    US-13 & Dashboards & En tant que RH/Admin, je consulte les KPI de sélection. & Moyenne \\
    \hline
    US-14 & Exports Excel & En tant que RH, j'exporte les suivis. & Moyenne \\
    \hline
    US-15 & Réclamations / notif. & Perspectives d'évolution. & Basse \\
    \hline
  \end{tabular}
\end{table}


\section{Planification des itérations --- répartition en sprints}

Compte tenu de la durée du PFA (environ deux mois) et de la nature progressive
de SCRUM, le travail a été découpé en sprints cohérents. Le Sprint~0 correspond
à la présente étude préliminaire. Les sprints suivants livrent des incréments
fonctionnels.

\begin{table}[htbp]
  \centering
  \small
  \caption{Répartition prévisionnelle des sprints}
  \label{tab:sprints-plan}
  \begin{tabular}{|c|p{3.6cm}|p{7.2cm}|}
    \hline
    \textbf{Sprint} & \textbf{Thème} & \textbf{Contenu principal} \\
    \hline
    0 & Étude préliminaire & Besoins, acteurs, backlog, architecture, techno \\
    \hline
    1 & Authentification \& profils & Login, rôles, profil, sécurité JWT \\
    \hline
    2 & Organisation Admin & Zones, agences, affectations RH, utilisateurs \\
    \hline
    3 & Offres et candidatures & CRUD offres, publication, dépôt CV, QCM, suivi \\
    \hline
    4 & Entretiens, décisions et pilotage & Meetings, embauche, dashboards, exports \\
    \hline
    5 & Module réclamations & Microservice CRUD réclamations (candidat / RH / Admin) \\
    \hline
  \end{tabular}
\end{table}

\begin{figure}[htbp]
  \centering
  \reportfig{img/gant.png}
  \caption{Diagramme de Gantt --- planification des sprints du projet}
  \label{fig:gantt-sprints}
\end{figure}

Cette répartition garantit d'abord un socle solide (identité, organisation),
puis le cœur du métier recrutement, et enfin le pilotage décisionnel.


\section{Conception architecturale}

Dans cette section, nous présentons l'architecture de la plateforme \textbf{DAAM}
afin de fournir une vue complète de ses composants fonctionnels.
L'architecture générale adopte un modèle à trois niveaux,
comprenant trois couches principales~: le front-end, le back-end sous forme
de microservices, et les bases de données, comme illustré dans la figure
ci-dessous.

\begin{figure}[htbp]
  \centering
  \reportfig{img/arch-three-tiers.png}
  \caption{Architecture générale à trois niveaux de la plateforme DAAM}
  \label{fig:arch-three-tiers}
\end{figure}

L'architecture de la plateforme DAAM repose sur une organisation en trois
couches principales afin d'assurer modularité, maintenabilité et
évolutivité.

\begin{itemize}
  \item \textbf{Le front-end (client)}~: conçu avec Angular, il gère
        l'interface utilisateur et permet aux acteurs (Administrateur, RH,
        Responsable de recrutement et Candidat) d'interagir avec le système
        via des formulaires, des listes, des tableaux de bord et un calendrier
        d'entretiens. Cette couche se charge de l'envoi des requêtes vers
        les services du back-end et de l'affichage des réponses.
  \item \textbf{Le back-end (serveur)}~: développé sous forme de
        microservices Spring Boot, il centralise la logique métier de
        l'application. Chaque service est dédié à un domaine
        spécifique (identité, recrutement, réclamations et notifications),
        ce qui favorise la scalabilité et facilite le déploiement
        indépendant de chaque composant.
  \item \textbf{La couche base de données}~: elle assure le stockage
        persistant des informations nécessaires au fonctionnement du
        système, notamment les comptes utilisateurs, l'organisation
        administrative (zones, entreprises), les données de recrutement
        (offres, candidatures, QCM, entretiens, embauches), ainsi que
        les réclamations et l'historique des notifications.
\end{itemize}

La communication entre ces couches s'effectue de manière synchrone via
HTTP/REST. Les ressources protégées s'appuient sur un jeton JWT transmis
par le client, garantissant un échange sécurisé et structuré entre
les composants.


\subsection{Conception et implémentation de l'architecture microservices}

L'architecture microservices repose sur une approche où l'application est
décomposée en plusieurs services autonomes, chacun ayant une
responsabilité métier spécifique. Cette conception permet une meilleure
scalabilité, modularité, maintenabilité et facilité de déploiement
indépendant de chaque composant.

Le projet DAAM est construit à l'aide de Spring Boot pour le back-end,
organisé en microservices, et d'un frontend Angular pour la couche client.
Les microservices retenus dans le périmètre du projet sont les suivants~:

\begin{itemize}
  \item \textbf{API Gateway}~: point d'entrée unique du système
        (port~8080). Il route les requêtes du frontend vers les
        microservices concernés et constitue la façade d'accès aux
        API REST.
  \item \textbf{Discovery Service (Eureka)}~: service de découverte qui
        permet aux microservices de s'enregistrer dynamiquement et de se
        localiser mutuellement sans configuration manuelle des adresses.
  \item \textbf{User Service}~: gère les opérations liées aux
        utilisateurs~: authentification JWT, création et gestion des
        comptes, profils, rôles (Admin, RH, Responsable de recrutement,
        Candidat) et droits d'accès.
  \item \textbf{Recruitment Service}~: porte le cœur métier du
        recrutement~: gestion des offres, candidatures et CV, QCM et
        scoring, suivi RH, planification des entretiens, décisions
        (retenu / non retenu / embauche), tableaux de bord et exports.
  \item \textbf{Reclamation \& Notification Service}~: un même
        microservice regroupant~:
        \begin{itemize}
          \item la \textbf{gestion des réclamations} (création, suivi,
                traitement et clôture) liées au parcours de recrutement~;
          \item l'\textbf{émission et la journalisation des notifications}
                métier (changement de statut, convocation à un entretien,
                rappel, confirmation d'embauche), avec possibilité
                d'évoluer vers plusieurs canaux (in-app, e-mail, etc.).
        \end{itemize}
        Ce choix mutualise le support utilisateurs et les alertes
        associées au sein d'un seul composant déployable.
\end{itemize}

À date, les services \textbf{User} et \textbf{Recruitment} sont opérationnels
dans l'application. Le microservice \textbf{Reclamation \& Notification}
est intégré à la conception architecturale et prévu comme extension du
système~; son implémentation complète constitue une perspective
d'évolution.

Chaque microservice dispose de sa propre base de données MySQL, et les
communications entre services s'effectuent principalement via des requêtes
HTTP RESTful.

Cette architecture microservices permet~:
\begin{itemize}
  \item une évolutivité horizontale, chaque service pouvant être
        répliqué et déployé indépendamment~;
  \item une meilleure résilience, car une panne dans un service
        n'impacte pas l'ensemble du système~;
  \item une gestion granulaire de la sécurité et des responsabilités
        métier~;
  \item un déploiement continu facilité par l'indépendance des
        composants.
\end{itemize}


\subsection{Architecture logique}

L'architecture logique du projet DAAM décrit les composants fonctionnels
et leurs interactions. Elle est organisée en plusieurs couches logiques~:

\begin{itemize}
  \item \textbf{Client}~: interface utilisateur Angular qui communique avec
        le backend via des appels HTTP. Elle gère la navigation par rôle
        et la conservation du jeton JWT côté navigateur.
  \item \textbf{API Gateway}~: couche de médiation assurant le routage
        des requêtes vers les microservices appropriés.
  \item \textbf{Microservices principaux}~:
        \begin{itemize}
          \item \textbf{User Service}~: authentification JWT, comptes,
                profils et permissions~;
          \item \textbf{Recruitment Service}~: offres, candidatures, QCM,
                entretiens, embauche, dashboards et exports~;
          \item \textbf{Reclamation \& Notification Service}~: dépôt /
                suivi des réclamations et diffusion des alertes métier
                au sein du même service.
        \end{itemize}
  \item \textbf{Services transversaux}~:
        \begin{itemize}
          \item \textbf{Discovery Service (Eureka)}~: localisation
                dynamique des microservices.
        \end{itemize}
  \item \textbf{Stockage de fichiers}~: conservation des CV et documents
        associés aux candidatures.
  \item \textbf{Bases de données}~: chaque microservice possède son
        propre schéma MySQL, assurant l'isolation des données
        (utilisateurs, recrutement, réclamations/notifications).
\end{itemize}

Les interactions se font principalement en REST synchrone, avec un
découplage fort grâce à l'indépendance des services et à l'usage
d'un service de découverte dynamique. Le service
\textbf{Reclamation \& Notification} pourra être sollicité par le
\textbf{Recruitment Service} lors d'événements clés du processus
(ex.~planification d'entretien, changement de statut), tout en exposant
ses propres API de support via le Gateway.

\begin{figure}[htbp]
  \centering
  \reportfig{img/arch-logic.png}
  \caption{Architecture logique de la plateforme DAAM}
  \label{fig:arch-logic}
\end{figure}


\subsection{Architecture physique (environnement d'exécution)}

L'architecture physique fait référence à l'infrastructure réelle qui
supporte le système logiciel. Dans le cadre du développement et des
tests, la plateforme DAAM s'exécute localement avec les composants
suivants~:

\begin{itemize}
  \item \textbf{Frontend Angular (port~4200)}~: accessible via le navigateur,
        il consomme les API REST exposées à travers la passerelle.
  \item \textbf{API Gateway (port~8080)}~: point d'entrée unique pour
        toutes les requêtes frontend~; il redirige les appels vers les
        microservices appropriés.
  \item \textbf{Eureka Server (port~8761)}~: registre central permettant la
        découverte dynamique des services.
  \item \textbf{Microservices Spring Boot} (communication synchrone)~:
        chaque service expose ses fonctionnalités via des API REST sur un
        port dédié~:
        \begin{itemize}
          \item User Service --- port~8089~;
          \item Recruitment Service --- port~8090~;
          \item Reclamation \& Notification Service --- port~8091 (prévu).
        \end{itemize}
  \item \textbf{Bases de données MySQL (port~3306)}~: chaque microservice
        dispose de son propre schéma, garantissant une séparation stricte
        des données.
  \item \textbf{Espace de stockage local}~: répertoire dédié aux
        fichiers CV téléchargés lors des candidatures.
\end{itemize}

Cette architecture assure~:
\begin{itemize}
  \item une haute disponibilité potentielle, chaque composant pouvant
        être répliqué~;
  \item une scalabilité horizontale, les microservices étant extensibles
        indépendamment~;
  \item un déploiement isolé, permettant des mises à jour ciblées
        sans interruption de l'ensemble du système~;
  \item une résilience accrue, limitant l'impact d'une panne à un seul
        service.
\end{itemize}

\begin{figure}[htbp]
  \centering
  \reportfig{img/arch-physical.png}
  \caption{Architecture physique de la plateforme DAAM}
  \label{fig:arch-physical}
\end{figure}

\section{Choix et exigences technologiques}

Les choix technologiques ont été retenus pour assurer une architecture robuste,
moderne et adaptée à une application web métier multi-rôles.


\subsection{Environnement matériel}

Le développement a été réalisé sur un poste de travail personnel présentant les
caractéristiques suivantes (à ajuster si nécessaire)~:
\begin{itemize}
  \item Processeur multi-cœurs récent~;
  \item Mémoire vive suffisante pour faire tourner frontend + microservices~;
  \item Stockage SSD~;
  \item Système d'exploitation Windows.
\end{itemize}


\subsection{Environnement logiciel}

\subsubsection{Frameworks et bibliothèques principales}

\begin{itemize}
  \item \textbf{Angular}~: construction de l'interface SPA et des espaces Admin /
        RH / Candidat.
  \item \textbf{ng-zorro-antd}~: composants UI (tables, formulaires, layout).
  \item \textbf{Spring Boot}~: développement des microservices REST.
  \item \textbf{Spring Security + JWT}~: sécurisation des API et gestion des
        autorités (rôles).
\end{itemize}


\subsubsection{Outils de développement}

\begin{itemize}
  \item \textbf{VS Code / Cursor}~: développement frontend et rédaction.
  \item \textbf{IntelliJ / IDE Java}~: développement backend.
  \item \textbf{Git}~: gestion de versions.
  \item \textbf{Postman} (ou équivalent)~: tests d'API.
\end{itemize}


\subsubsection{Langages}

\begin{itemize}
  \item \textbf{TypeScript / HTML / CSS}~: frontend.
  \item \textbf{Java}~: backend Spring Boot.
  \item \textbf{SQL}~: persistance relationnelle.
\end{itemize}


\section{Conclusion}

Le Sprint~0 a permis de poser les fondations du projet~: acteurs clairement
identifiés, besoins fonctionnels et non fonctionnels formalisés, Product Backlog
priorisé, découpage en sprints, et architecture technique définie.

Cette étude préliminaire garantit que les itérations suivantes s'appuient sur
une vision partagée du produit et sur une trajectoire de livraison réaliste.
Le chapitre suivant présente le Sprint~1, consacré à l'authentification, aux
profils et à la gestion de base des utilisateurs.

\clearpage
